\chapter*{Abstrak}

\makeatletter
    \vspace{1cm}
    \begin{center}
        \noindent\textbf{\@judul}
    \end{center}

    \vspace{1cm}
    \begin{table}[htbp]
        \begin{tabular}{lcp{0.49\paperwidth}}
            Nama Mahasiswa / NRP & : & \@nama~/ \@NRP\\
            Departemen & : & \@departemen, Fakultas ~{\@fakultas} ITS\\
            Calon Dosen Pembimbing & : & \pbb@satu
        \end{tabular}
    \end{table}

    \vspace{1cm}
    \noindent
    \textbf{Abstrak}

    % ============================
    % Teks abstrak mulai dari sini
    % ============================
    Permasalahan \textit{reward balancing} dalam desain ekonomi gim merupakan salah satu jantung keberhasilan dari suatu gim tersebut. Sistem \textit{reward} yang tidak seimbang menyebabkan menurunnya ketertarikan pemain, percepatan progres permainan yang tidak wajar, serta dapat memengaruhi keberlanjutan ekonomi in-gim secara keseluruhan. Beberapa penelitian sebelumnya telah membahas mengenai sistem ekonomi gim secara umum, namun pendekatan tersebut masih berfokus pada skala makro dan belum mempertimbangkan dinamika perilaku pemain. Penelitian ini bertujuan untuk mengembangkan optimasi sistem reward yang mampu mempertimbangkan variasi perilaku pemain dengan menggunakan algoritma evolusioner diintegrasikan dengan simulasi agen. Setiap agen digunakan untuk mewakili dinamika perilaku pemain yang berinteraksi dalam lingkungan gim. Algoritma evolusioner digunakan untuk mendapatkan kombinasi optimal parameter \textit{reward} yang menghasilkan sistem paling seimbang. Dengan menekankan pendekatan komputasional-matematis, penelitian diharapkan dapat menghasilkan konfigurasi sistem \textit{reward} yang seimbang dan berkelanjutan pada berbagai variasi skenario pemain sekaligus memberikan kontribusi ilmiah bagi perancangan sistem ekonomi gim yang dinamis dan tahan terhadap perubahan.
    
    % Untuk kata kunci
    \katakunci{Optimasi, Ekonomi, Reward}
\makeatother